\lysh{34814}{2}
\begin{alist}
    \item Από τον πίνακα βλέπουμε ότι η τελευταία αθροιστική συχνότητα είναι $N_5=40$ άρα το μέγεθος του δείγματος είναι $ν=40$.
    \item
    $N_1=5$ άρα και $ν_1=5$
    \par $N_2=15$ και επειδή $N_2=N_1+ν_2 \Leftrightarrow 15=5+ν_2 \Leftrightarrow ν_2=10$
    \par $N_3=20$ και επειδή $N_3=N_2+ν_3 \Leftrightarrow 20=15+ν_3 \Leftrightarrow ν_3=5$
    \par $N_4=35$ και επειδή $N_4=N_3+ν_4 \Leftrightarrow 35=20+ν_4 \Leftrightarrow ν_4=15$
    \par $N_5=40$ και επειδή $N_5=N_4+ν_5 \Leftrightarrow 40=35+ν_5 \Leftrightarrow ν_5=5$
    \par Άρα ο πίνακας συμπληρώνεται ως εξής:
    
    \begin{center}
\begin{mytblr}{}
        Κλάσεις $[ - )$ & Συχνότητα $\nu_i$ & {Αθροιστική\\Συχνότητα $N_i$} & {Κεντρική\\τιμή $x_i$} &  $x_i \cdot \nu_i$ \\
        0-2 & \textcolor{maincolor}{5} & 5 & 1 & \textcolor{maincolor}{5} \\
        2-4 & \textcolor{maincolor}{10} & 15 & 3 & \textcolor{maincolor}{30} \\
        4-6 & \textcolor{maincolor}{5} & 20 & 5 & \textcolor{maincolor}{25} \\
        6-8 & \textcolor{maincolor}{15} & 35 & 7 & \textcolor{maincolor}{105} \\
        8-10 & \textcolor{maincolor}{5} & 40 & 9 & \textcolor{maincolor}{45} \\
        
        ΣΥΝΟΛΟ & $40$ & \SetCell{bg=gray!30} & \SetCell{bg=gray!30} & \textcolor{maincolor}{\bfseries 210} \\
    \end{mytblr}
\end{center}

    \item Γνωρίζουμε ότι η μέση τιμή ισούται με:
    \[
    \bar{x} = \frac{x_1 \cdot \nu_1 + x_2 \cdot \nu_2 + x_3 \cdot \nu_3 + x_4 \cdot \nu_4 + x_5 \cdot \nu_5}{40} = \frac{210}{40} = 5,25
    \]
\end{alist}
\vspace{6mm}
\noindent\rule{\textwidth}{0.4pt}
\vspace{6mm}

